\chapter{绪论}

\section{研究背景与意义}

无人机是依靠遥控设备或机载程序执行飞行任务的不载人飞行器，具有部署速度快、运行成本低和环境适应性强等特点。随着飞行平台、小型动力系统、导航感知、通信链路和自主控制技术的持续发展，无人机执行物流运输任务的工程条件逐步成熟\cite{GDJX202501002,JTJS202503011}。相较于依赖地面路网的传统配送方式，低空物流能够利用立体空域完成点到点运输，在应急物资投送、医疗样本转运、山区补给和城市末端配送等场景中显著缩短运输时间，并在道路通达性较差、配送节点分散或时效要求严格的任务中扩大运输可达范围，减轻对地面路网条件的依赖。图~\ref{fig:uav_transportation} 展示了无人机在山地物资运输和城市物流等不同环境中的典型应用形态。随着低空物流由示范运行逐步转向常态化服务，无人机系统已不再仅需完成简单到达任务，而需要在受限空域内统筹建筑群、禁飞区域、复杂地形和起降条件等多种约束，路径规划、轨迹生成、任务决策与飞行控制之间的衔接也更为紧密。

\begin{figure}[htbp]
    \centering
    % 定义第一个子图容器，宽度占正文宽度的 45%
    \begin{minipage}[t]{0.45\textwidth}
        \centering
        % 强制图像宽度充满容器（即 0.45\textwidth），高度固定为 4cm
        % 移除 keepaspectratio 参数以实现绝对尺寸对齐
        \includegraphics[width=\textwidth, height=4cm]{pic/part0/无人机1.png}
        \par\vspace{0.2cm} 
        \small （a）无人机山地物资运输
    \end{minipage}%
    \hspace{0.05\textwidth} % 设定子图间的固定水平间距
    \begin{minipage}[t]{0.45\textwidth}
        \centering
        \includegraphics[width=\textwidth, height=4cm]{pic/part0/无人机2.png}
        \par\vspace{0.2cm}
        \small （b）无人机城市物流
    \end{minipage}
    \caption{无人机在不同环境执行物流运输任务}
    \label{fig:uav_transportation}
\end{figure}

当前低空物流正由单架无人机、单批次配送逐步发展为面向多节点、多频次和多类型货物的持续运行模式。轻小件高频配送强调快速响应和系统吞吐能力，重载运输任务则更加依赖多机协同、精确控制和稳定执行；与此同时，单架无人机在载荷、续航和服务半径上的限制日益凸显，当配送范围扩大、任务节点增多或货物重量超过单机承载上限时，仅依赖单机往往难以同时满足效率、安全与经济性要求。多无人机系统通过任务分工、空间协作和能力互补，能够提升单位时间内的配送能力、扩展任务覆盖范围，并为单机难以完成的重载吊运提供可执行方案，因此已成为高频、多点和重载低空物流任务中的重要组织形态。图~\ref{fig:mutiuav} 给出了多无人机物流系统在不同任务中的典型组织形式。

\begin{figure}[htbp]
    \centering
    % 定义第一个子图容器，宽度占正文宽度的 45%
    \begin{minipage}[t]{0.45\textwidth}
        \centering
        % 强制图像宽度充满容器（即 0.45\textwidth），高度固定为 4cm
        % 移除 keepaspectratio 参数以实现绝对尺寸对齐
        \includegraphics[width=\textwidth, height=4cm]{pic/part0/物流系统1.png}
        \par\vspace{0.2cm} 
        \small （a）无人机山地物资运输
    \end{minipage}%
    \hspace{0.05\textwidth} % 设定子图间的固定水平间距
    \begin{minipage}[t]{0.45\textwidth}
        \centering
        \includegraphics[width=\textwidth, height=4cm]{pic/part0/物流系统2.png}
        \par\vspace{0.2cm}
        \small （b）无人机城市物流
    \end{minipage}
    \caption{多无人机物流系统的典型任务}
    \label{fig:mutiuav}
\end{figure}

多无人机系统在带来能力提升的同时，也显著增加了任务组织复杂度。多架无人机需要在共享空域内协调航迹与时序、保持整体任务目标一致，编队飞行、群体避让、协同吊运和分布式执行等过程都会引入机间约束和状态耦合，使低空物流系统呈现出明显的分层结构与跨层耦合特征。结合低空物流任务链路，复杂环境下的路径规划、多无人机飞行过程中的轨迹规划与安全避碰，以及协同运输过程中的建模与控制，构成了关系任务可达性、空域安全与运输稳定性的核心问题。

在任务起始阶段，全局路径规划直接影响无人机能否在高密度建筑、受限空域和多样飞行约束并存的城市低空环境中高效到达目标区域，若仅关注几何最短距离而忽略飞行能耗特性，所得路径往往难以充分发挥平台续航能力；而在获得全局路径后，多无人机还需在共享空域内同时满足障碍物避让、机间安全间隔保持和任务时序协调等要求，定位误差、环境不确定性和多机相互影响又会进一步放大实时轨迹生成与安全避碰的难度。因此，如何在复杂三维环境和不确定条件下统筹路径可行性、能耗效率与多机安全，是低空物流运行中的关键研究问题。

当任务进一步进入重载运输阶段，多无人机协同吊运虽然能够显著扩展系统能力，但无人机、绳索与负载之间的相互作用会带来强非线性、欠驱动和多体耦合特征，控制输入、绳索张力与负载运动之间存在显著联动，对建模精度和控制稳定性提出更高要求。因此，如何构建适用于复杂耦合系统的建模与控制方法，是多无人机运输走向实用化必须解决的问题。围绕上述问题，开展能耗约束下基于改进关键可视图的路径规划、考虑位置不确定性的多无人机轨迹规划以及基于强化学习的多无人机协同运输控制研究，在统一的低空物流任务链路中形成从规划到控制的系统化研究框架。

\section{国内外研究现状}

围绕低空物流中的多无人机作业需求，国内外学者已在路径规划、多无人机轨迹规划以及协同运输控制等方向开展了大量研究。现有研究已经形成了较为丰富的方法积累，但面向低空物流任务时，考虑多种约束的路径规划、存在位置误差下多无人机协同轨迹规划，以及协同运输控制中的强耦合动力学问题仍需进一步研究，以下分别对上述三个方向的研究现状进行归纳分析。

\subsection{无人机路径规划研究现状}

现有无人机路径规划研究可按照环境表示方式和路径生成机理分为单元分解法、路线图法、人工势场法和启发式优化方法\cite{HKGC202302003,HKJJ202204030}。前两类方法侧重于环境连通结构的构造与搜索，人工势场法强调基于局部势场的反应式规划，启发式优化方法则通常以离散航路点序列为优化对象，通过迭代搜索获得满足约束的可行路径。

单元分解法通过将环境空间划分为若干互不重叠的单元，并基于单元之间的邻接关系生成可行路径。根据环境划分方式的不同，现有研究中的单元分解法可进一步分为近似单元分解、精确单元分解和自适应单元分解三类。近似单元分解通常采用规则网格对空间进行离散表示，方法实现简单，便于在二维或三维空间中构建统一的搜索结构，但其环境表达精度受网格分辨率影响明显。Han等提出了一组面向复杂室内环境的无人机网格优化路径规划算法，先利用 GeoSOT-3D 对室内三维空域进行网格化建模，再分别设计面向常规环境的 GO-APP 和面向复杂“死区”环境的 GO-LBPP，以降低计算复杂度、加快收敛并改善路径质量\cite{han2022grid}。Chiciudean提出了面向无人机导航的三维网格路径搜索方法，先对三维环境进行预计算离散表示，再在该三维网格中搜索无碰且接近最优的飞行轨迹\cite{chiciudean2022pathfinding}。精确单元分解依据障碍物几何边界直接将环境划分为若干多边形区域，典型方式包括梯形分解和往复式分解，该类方法能够更贴合障碍边界表达环境连通关系。Choset和Pignon提出了往复式单元分解方法，通过在环境连通性发生变化的位置构造精确分解单元，使每个单元内部均可通过简单往返运动完成覆盖，并在单元串联基础上形成全局覆盖路径\cite{choset1998coverage}。Coombes提出了面向固定翼无人机测绘任务的最优多边形分解方法，将空间分解与风场影响、飞行时间优化直接结合起来\cite{coombes2018optimal}。自适应单元分解则根据障碍物分布情况对局部空间递归细分，仅在必要区域提高分辨率，典型代表包括四叉树和八叉树分解方法。Finkel提出了四叉树数据结构，为路径规划中“仅在必要区域递归细分”的环境表达方法提供了基础\cite{finkel1974quad}。Kitamura提出了基于八叉树和人工势场的三维动态环境路径规划方法，利用八叉树表示三维动态空间，再结合人工势场生成无碰路径\cite{kitamura1995path}。该类方法具有建模直观、结构清晰和易于实现等特点，适合将复杂环境转换为可搜索的离散空间，但对空间离散粒度较为敏感：分解过细会导致节点规模迅速增加，抬高搜索代价；分解过粗又容易造成障碍边界表达失真和路径质量下降。在三维复杂环境中，随着障碍数量增多和环境尺度扩大，单元数量及其邻接关系会快速膨胀，从而使规划效率受到明显影响。因此，如何在环境表达精度与求解复杂度之间取得平衡，仍是单元分解法需要解决的问题。

路线图法通过从自由空间中提取具有代表性的连通骨架或图结构，将连续空间路径规划问题转化为图上的路径搜索问题。根据图结构构造方式的不同，路线图法可进一步分为可视图法、Voronoi图法、概率路图法和快速随机扩展树法。可视图法通过连接障碍物顶点之间的可见边构造可视图（Visibility Graph，VG）模型，能够较直接地反映自由空间中的可通行关系，通常具有较好的路径最短性，但生成路径往往贴近障碍物边界。Lozano-Pérez提出了面向多面体障碍环境的无碰路径规划算法，通过对障碍物进行几何膨胀，将碰撞约束转化为参考点的禁入区域，再在障碍顶点之间构造安全连通网络，奠定了可视图法的经典几何框架\cite{lozano1979algorithm}。Blasi提出了基于关键可视图的路径规划与避碰方法，在关键可视图（Essential Visibility Graph，EVG）上搜索最小代价分段线性路径，并通过图更新实现动态障碍条件下的重规划\cite{blasi2020path}。随后，Blasi进一步提出分层关键可视图（Layered Essential Visibility Graph，LEVG），将二维EVG扩展到三维约束环境\cite{blasi2023uav}。Hu等提出了三维分层可视图航路网络模型，并在此基础上开展多目标路径规划，同时考虑路径长度、风险与运行约束\cite{hu2025research}。Voronoi图法通过构造与障碍物边界等距的骨架路径，提高了路径对障碍物的安全间隔，更适合强调安全裕度的场景。Beard等在多无人机协同目标分配与拦截问题中引入了基于威胁的Voronoi图，通过在图边上引入威胁代价和长度代价，引导无人机沿远离威胁的安全骨架飞行\cite{beard2002coordinated}。Dong等提出了Voronoi--A-star算法（A-star，A*）融合算法，用于复杂地形中的UAV路径规划，先根据障碍密度对高程地形进行自适应分层，再在各层内构造稀疏节点网络，并结合A*完成局部补充搜索\cite{dong2025voronoi}。概率路图法（Probabilistic Roadmap，PRM）通过在自由空间中随机布设连接节点并建立邻接关系，形成具有概率完备性的连通图结构，对复杂空间具有较强适应性\cite{kavraki1996probabilistic}。Yan等提出了面向复杂三维环境的PRM路径规划算法，先用八叉树划分环境，再通过在包围盒区域内采样改进路图节点分布，最后在所构建的路图上搜索可行路径\cite{yan2013path}。Jin等提出了面向室内安全飞行的改进PRM方法，通过点云投影和多层栅格表示改进传统PRM的路网构建、搜索顺序与路径后处理过程\cite{jin2023improved}。快速随机扩展树法（Rapidly-exploring Random Tree，RRT）通过在状态空间中随机采样并逐步扩展树结构实现路径搜索，其特点是不预先构建完整图模型，而是从起点出发边采样边生长搜索树，因此能够较快探索高维连续空间。LaValle提出了RRT，将路径规划问题转化为一棵从起点不断向随机采样状态扩展的搜索树，使算法能够快速探索高维状态空间\cite{lavalle1998rapidly}。Karaman提出了RRT*和PRM*，并证明其具有渐近最优性\cite{karaman2011sampling}。Ramana等研究了固定翼UAV在城市场景中的RRT运动规划问题，重点处理短时间窗内近优可飞路径生成\cite{ramana2016motion}。路线图法强调对自由空间主要连通结构的抽象表达，通常能够以较少的节点和边描述环境中的关键通道，因此在复杂环境中具有较高的搜索效率和较好的可解释性，但其效果同样依赖前期图结构构建质量。对于障碍密集或空间结构复杂的三维环境，图构建过程本身可能带来较高的计算成本；当任务中需要进一步引入风险、通行规则或多种任务代价时，还需要在图模型层面对边权和连通关系进行额外设计。因此，如何在保持图结构紧凑性的同时增强其对复杂任务约束的表达能力，仍是路线图法研究中的重要问题。

人工势场法通过构造目标点对无人机的吸引势场以及障碍物对无人机的排斥势场，引导无人机沿合成势场方向连续运动，从而实现路径规划。该类方法不依赖显式全局图模型，形式简洁，计算量较小，具备较好的实时响应能力，适用于局部避障和动态环境中的快速反应式规划。由于势场能够随障碍物状态实时变化而更新，人工势场法在在线路径调整和局部导航中具有一定应用价值。Jayaweera等提出了动态人工势场无人机路径规划方法，用于跟踪地面运动目标，将目标运动信息直接纳入势场构造过程\cite{jayaweera2020dynamic}。Hao等提出了改进人工势场路径规划方法，通过引入碰撞风险评估机制、虚拟子目标以及改进势场函数，缓解了传统人工势场法中的局部极小、目标不可达和绕障不合理等问题\cite{hao2023uav}。Shao等提出了基于强化学习的三维自适应人工势场算法，在三维状态空间中利用学习策略动态调整排斥势场增益，并对引力函数、目标附近势场和奖励机制进行联合设计\cite{shao2025adaptive}。然而，人工势场法的局限同样明显：由于路径生成依赖局部势场分布，规划结果容易受到局部极小值、局部振荡以及引力与斥力失衡等问题影响，进而削弱路径搜索的稳定性与全局可达性；同时，该方法对局部环境感知和势场信息完整性依赖较强，若环境信息不足或势场表达不准确，则会显著影响规划效果。因此，在障碍密集、通道狭窄或环境复杂程度较高的场景中，人工势场法通常需要与其他全局规划或引导策略结合使用，以提高路径质量和算法鲁棒性。

启发式优化方法通常将路径表示为一组有序离散航路点，并通过遗传算法、粒子群算法、蚁群算法等启发式策略，对航路点位置、数量或连接顺序进行迭代调整，以获得满足约束条件的较优路径\cite{KZYC202205002}。相较于前述环境表示方法，启发式优化方法具有更强的目标函数建模灵活性，能够将路径长度、威胁规避、约束惩罚和任务代价等因素统一纳入优化框架，因此在复杂约束和多目标规划问题中具有较强适应性。Niu等提出了改进的连续域蚁群优化方法，在连续航路点搜索框架下引入改进状态转移概率、随机游走策略和自适应航点修复机制，以提升三维UAV路径规划的收敛速度、路径质量与可行性，并进一步与动态窗口法结合用于部分未知场景中的静态和动态障碍避让\cite{niu2025path}。Gutierrez-Martinez等提出了面向三维静态和动态环境的四旋翼无人机遗传算法路径规划方法，通过重新设计遗传操作、较高变异准则、种群重置机制以及目标点嵌入策略，使遗传算法能够在实时任务中生成优化轨迹\cite{gutierrezmartinez2025genetic}。Qi等针对多无人机路径规划中的多类能耗约束问题，建立了考虑平飞、加速、减速、爬升和转向能耗的模型，并提出改进的蜂觅学习粒子群优化算法，提高了多机路径规划的求解效率与稳定性\cite{qi2025multi}。然而，启发式优化方法的性能高度依赖优化算法本身。对于高维搜索空间或复杂障碍环境，相关方法容易出现收敛速度下降、参数敏感性增强以及结果波动较大的问题；若离散航路点数量设置过多，会显著扩大优化规模，增加求解负担；若设置过少，则可能限制路径表达能力。因此，如何在优化灵活性、求解效率与结果稳定性之间取得平衡，仍是启发式优化方法研究中的关键问题。

总体来看，现有无人机路径规划研究在环境建模、连通路径搜索、局部避障和路径优化等方面已经形成了较为丰富的方法体系。面向低空物流任务时，路径规划不仅关系到飞行可达性，也直接影响续航利用率和任务经济性。现有研究对飞行能耗的考虑仍偏粗略，在能耗建模精度、复杂环境适应性与规划效率之间尚缺少更有针对性的权衡，围绕能耗约束的路径规划方法仍需进一步研究。

\subsection{多无人机轨迹规划研究现状}

现有多无人机轨迹规划研究大致可分为集中式协同轨迹规划、分布式协同轨迹规划以及面向编队保持的轨迹规划三类，近年来强化学习等数据驱动方法也逐渐进入无人机集群航迹规划研究\cite{KJDZ202506001}。集中式方法侧重利用全局环境信息和多无人机状态进行统一建模与联合求解，分布式方法强调各无人机基于局部信息或有限通信完成轨迹生成与协同决策，编队轨迹规划方法则更加关注多无人机作为整体在飞行过程中的队形保持、结构组织与环境适应问题。

集中式协同轨迹规划方法通过中心决策单元统一获取环境信息与无人机状态，对多机轨迹进行整体规划与协调求解。根据具体方法机理的不同，现有研究可进一步分为基于集中式联合优化的方法、基于集中式模型预测控制的方法以及基于集中式冲突检测与冲突消解的方法。基于集中式联合优化的方法通常将多架无人机的轨迹、时序和安全约束统一写入同一优化模型，通过联合求解获得整体协调轨迹。Zhang等面向多无人机货运系统，联合考虑蜂窝通信约束、任务分配和轨迹优化问题，通过统一优化框架协调多机执行过程\cite{zhang2025cargo}。Yang等提出了双层协同规划算法，在四维约束条件下同时处理航迹生成与时序协调问题\cite{yang2025twolayer}。基于集中式模型预测控制的方法则在统一预测时域内滚动生成满足动态约束和安全约束的协同轨迹。Wang等将多无人机协同跟踪问题建模为集中式模型预测控制问题，在城市环境中通过联合优化提高目标可见性\cite{wang2022visibility}。Chen等面向复杂环境中的非合作目标跟踪，构建多无人机协同轨迹规划模型，并通过非线性模型预测控制实现统一滚动优化\cite{chen2025nmpc}。基于集中式冲突检测与冲突消解的方法通常先识别多机飞行过程中的潜在冲突，再通过统一调整路径、速度或时序实现整体冲突消解。Feng等针对狭窄环境下的多无人机协同路径规划问题，采用分层冲突搜索与低层轨迹生成相结合的方式处理多机冲突\cite{feng2024cramped}。Lu等结合改进跳点搜索与冲突检测机制，实现了面向三维环境的多无人机无冲突路径规划\cite{lu2025conflict}。集中式方法能够较好地处理机间耦合关系和整体任务约束，但其计算复杂度与通信负担通常随无人机数量增加而迅速上升，在大规模系统和复杂环境中的实时应用受到限制。

分布式协同轨迹规划方法强调各无人机基于局部感知信息、邻域信息或有限通信信息自主完成轨迹生成，并通过协同机制实现整体安全飞行与任务协调。根据具体方法机理的不同，现有研究可进一步分为基于分布式优化的方法、基于模型预测控制的分布式方法、基于速度障碍与相对运动几何的方法以及基于运动原语（Motion Primitives，MP）的方法。基于分布式优化的方法通常将整体规划问题分解为多个局部子问题，并通过局部求解与信息交互逐步协调一致。Liu等研究了有向通信网络下的多无人机分布式轨迹优化问题，通过主从分解框架实现分布式协调求解\cite{liu2021distributed}。Duan等面向动态且信息丰富的环境，提出分布式协同路径规划方法，以增强多机系统在动态场景中的适应能力\cite{duan2025distributed}。基于分布式模型预测控制（Distributed Model Predictive Control，DMPC）的方法在局部预测时域内滚动生成满足约束的可行轨迹。Yang等提出分布式模型预测编队控制方法，在障碍环境中实现多机协同飞行与编队保持\cite{yang2025dmpc}。Gräfe等面向微小型无人机集群提出DMPC-Swarm框架，在有限通信和消息丢失条件下实现分布式滚动控制\cite{grafe2025dmpc}。基于速度障碍与相对运动几何的方法通过分析无人机之间及其与障碍物之间的相对速度关系构造安全机动策略，具有较好的实时性。Xian等提出三维速度障碍多无人机避碰策略，通过构建速度禁区实现多机安全机动\cite{xian2025vo}。Tao等结合碰撞锥和预警判据设计反应式避碰策略，提高了相对运动几何方法在在线避障中的有效性\cite{tao2023rcatrs}。基于 MP 的方法通过预先构造满足动力学约束的局部可行轨迹片段，并在线选择或拼接候选轨迹，以提高复杂环境中的快速响应能力。Manzoni等提出基于 MP 的无人机障碍环境轨迹规划方法，通过构造满足机动约束的原语集合并结合在线搜索，提高了复杂环境中的实时规划效率\cite{manzoni2024mprims}。Zhang等面向大规模入侵无人机集群拦截问题，利用 MP 建立实时轨迹规划框架，并分析了在群体目标场景下的拦截效率\cite{zhang2024swarmprimitive}。分布式方法在可扩展性和实时性方面具有明显优势，但在复杂障碍环境中仍容易受到局部最优、反复冲突以及全局协调不足等问题影响。

基于编队的轨迹规划方法将多无人机系统视为具有特定空间结构的整体，重点研究在保持编队形状基础上的轨迹生成问题。根据具体组织方式的不同，现有研究主要包括基于领导者—跟随者的编队轨迹规划方法、基于虚拟结构的编队轨迹规划方法以及基于行为规则的编队轨迹规划方法。基于领导者—跟随者的方法通常由领导者生成参考轨迹，其余无人机依据相对位置关系进行跟踪。Sheng等提出改进人工势场编队保持方法，在领导者—跟随者框架下提升了障碍环境中的编队保持能力\cite{sheng2023leader}。Zhu等利用多参与者Stackelberg--Nash博弈方法研究固定翼无人机领导者—跟随者编队重构控制问题\cite{zhu2025leader}。基于虚拟结构的方法将整个编队视作一个具有整体几何形状的虚拟体，通过规划整体轨迹实现编队运动。Guo等围绕虚拟结构方法开展多无人机编队与半实物仿真研究，验证了基于整体几何结构组织编队飞行的有效性\cite{guo2023virtual}。Wang等进一步将虚拟结构与深度强化学习相结合，用于提高复杂环境中的编队控制能力\cite{wang2025virtual}。基于行为规则的方法通过设计吸引、排斥和对齐等局部交互规则维持编队结构并实现协同运动。Yue等面向固定翼无人机群提出考虑安全约束的行为型协同控制方法，通过虚拟管道机制实现群体组织与安全飞行\cite{yue2025behavior}。Liu等结合深度强化学习与仿射变换方法，实现了多无人机柔性编队控制与轨迹调整\cite{liu2025flexible}。该类方法能够较好地保持群体有序性并支持整体机动，但当编队需要穿越狭窄区域或在障碍密集环境中飞行时，队形保持与安全通行之间往往存在明显冲突，编队结构的灵活调整和整体安全约束处理仍然是研究中的关键问题。

总体来看，现有多无人机轨迹规划研究已经覆盖集中式、分布式和编队三类主要技术路线。面向障碍物密集且存在无人机位置不确定性的低空物流任务时，集中式方法在大规模条件下实时性受限，分布式方法容易出现全局协调不足的问题，基于编队的方法则难以同时兼顾队形保持、环境适应与整体安全。尤其在存在定位误差和环境不确定性的条件下，如何将编队整体碰撞风险建模、实时轨迹生成与安全约束统一纳入规划框架，仍是多无人机轨迹规划需要解决的关键问题。

\subsection{多无人机协同运输控制研究现状}

现有多无人机协同运输控制研究主要沿三条思路展开，包括基于系统动力学建模的方法、基于抗扰动编队控制的方法以及基于强化学习的方法\cite{XXYK202101006,KZYC202504002}。第一类方法从系统动力学出发，对无人机、负载及连接约束进行统一建模，并在此基础上设计控制器；第二类方法将负载影响等效为外部扰动、附加力或编队约束，借助已有编队控制框架完成运输任务；第三类方法则通过环境交互学习控制策略，用于处理复杂非线性和多约束耦合条件下的协同运输控制问题。

基于系统动力学建模的方法通常从多无人机、负载以及绳索或刚性连接构成的整体系统出发，对系统运动学与动力学关系进行统一建模，并在此基础上设计几何控制、反馈线性化、逆动力学分配、鲁棒控制或自适应控制等控制律。Lee建立了四旋翼运输缆绳悬挂刚体负载的几何动力学模型，并设计了相应的几何控制方法，为协同运输系统的非线性建模与稳定控制提供了基础\cite{lee2018geometric}。Klausen等针对未知负载质量和环境扰动条件下的多旋翼协同运输问题，提出了合作控制方法，提高了模型驱动控制在不确定工况下的适用性\cite{klausen2020cooperative}。Geng和Langelaan提出了负载主导的分层控制框架，由负载层生成系统所需合力与合力矩，再分配至各飞行器执行控制\cite{geng2020cooperative}。Chen等针对双四旋翼协同运输柔性负载问题，建立了柔性梁负载动力学模型并设计层级控制结构，将协同运输对象由刚体扩展到柔性负载\cite{chen2021cooperative}。Chai等提出了能量型非线性自适应控制方法，提高了系统对负载摆动和外部扰动的抑制能力\cite{chai2022energy}。该类方法能够较准确地描述系统耦合关系，在轨迹跟踪、姿态调节和张力分配等方面具有较强的理论基础，但其对模型精度、参数先验和状态测量条件依赖较强，且随着系统规模扩大，建模复杂度与在线求解负担也会明显增加。

基于抗扰动编队控制的方法通常将运输负载对无人机系统的影响等效为外部扰动、附加力或编队保持过程中的等效约束，并利用已有多无人机编队控制、分布式协调控制或一致性控制框架实现运输任务。Doakhan等针对多四旋翼协同运输中的不确定性问题，设计了实时编队控制结构，将负载轨迹跟踪与编队保持统一纳入控制框架\cite{doakhan2023realtime}。Doakhan等进一步提出鲁棒自适应编队控制方法，利用自适应滑模结构处理参数不确定性与外部扰动\cite{doakhan2023robust}。Su等将协同运输问题转化为有界扰动条件下的编队跟踪问题，提出鲁棒自适应编队控制方法，并通过参数修正机制减小不确定性带来的性能退化\cite{su2023robust}。Ccari和Yanyachi将神经网络逼近、自适应律与编队控制结合，用于两架欠驱动四旋翼的协同负载运输控制\cite{ccari2023novel}。Li等考虑动态通信拓扑切换条件下的四机协同运输问题，结合有限时间扩张状态观测器与一致性控制方法，实现了对负载扰动的快速估计与编队保持控制\cite{li2026formation}。该类方法便于利用已有编队控制成果，在工程实现和系统扩展方面具有优势，但当负载质量较大、摆动显著或多机之间存在明显受力重分配时，扰动等效处理难以准确反映绳索张力变化、负载姿态演化和多机协同受力调整等关键机制，因此对强耦合运输过程的刻画能力有限。

基于强化学习的方法通过构建任务环境、状态空间、动作空间与奖励机制，使控制器能够通过交互数据直接学习协同运输策略。Faust等将强化学习用于悬挂载荷无人机的运输控制，实现了载荷摆振抑制与轨迹生成\cite{faust2017automated}。Belkhale等提出模型驱动元强化学习方法，使飞行器能够在负载参数变化后快速适应运输控制任务\cite{belkhale2021model}。Lin等针对双四旋翼协同负载运输问题，采用集中式强化学习方法实现位置控制与系统稳定\cite{lin2024payload}。Chen等提出深度多智能体强化学习框架，使多无人机能够自主导航至负载抓取点，为抓取、起吊与运输一体化协同控制提供了基础\cite{chen2023deep}。Chen等进一步提出面向双机负载运输的元强化学习方法，将在线适应、动作修正与碰撞感知规划结合起来，提高了系统在风扰和构型变化条件下的运输能力\cite{chen2025meta}。该类方法能够在较少依赖精确模型反解的条件下处理复杂非线性、高维状态和多约束耦合问题，但当前研究仍受到训练效率、样本成本、奖励设计、收敛稳定性和高保真验证条件等因素制约，直接面向多无人机协同负载闭环控制的研究仍然较少。

总体来看，现有多无人机协同运输控制研究已经形成了基于模型、基于编队控制框架扩展以及基于强化学习的多类技术路线。基于系统动力学建模的方法能够较准确地描述多无人机、负载及连接约束之间的耦合关系，但模型构建复杂，对参数精度和在线求解能力要求较高。基于抗扰动编队控制的方法便于利用已有编队控制框架实现协同运输，但对强耦合运输过程中的张力变化、负载摆动和受力重分配等关键机制刻画不足。基于强化学习的方法能够在较少依赖精确模型反解的条件下处理复杂非线性和多约束耦合问题，但仍面临训练效率、策略稳定性和泛化能力等方面的限制。面向低空物流任务中的多无人机协同运输问题，复杂耦合关系、控制实时性和任务适应性仍需统一考虑。


\section{研究内容和技术路线}

技术路线从低空物流任务需求出发，在全局规划层采用能耗约束下基于改进关键可视图的路径规划方法，为无人机生成复杂环境中的可行任务路径；在局部规划层采用考虑位置不确定性的多无人机轨迹规划方法，为多无人机编队生成满足安全约束的飞行轨迹；在运输执行层采用基于强化学习的多无人机协同运输控制方法，并利用基于 JAX/MJX 的并行化仿真环境提升协同运输控制的训练效率；三部分工作分别对应全局航路生成、局部安全通行与协同运输控制，形成由规划到控制逐层细化的研究链路。围绕上述技术路线，研究内容分为以下三部分，具体技术路线如图~\ref{fig:overall} 所示。

\begin{figure}[!h]
  \centering
  \includegraphics[width=1\textwidth]{pic/research_content_overview.drawio.png}
  \caption{研究内容示意图}
  \label{fig:overall}
\end{figure}

\subsection{能耗约束下基于可视图的路径规划方法}

针对低空物流任务中复杂三维环境下的全局路径生成问题，提出能耗约束下基于可视图的路径规划方法。该部分工作以关键可视图方法为基础，在三维城市环境中对障碍物进行筛选和分层表达，并结合无人机动力学能耗建模构造路径代价函数。针对传统可视图在三维环境中节点冗余较多、可视性检查开销较大的问题，通过关键障碍筛选、候选节点生成和可视性判定方法改进图结构构建过程，再结合同时考虑路径长度与能耗约束的搜索策略，获得满足低空物流任务需求的全局路径。该部分研究为后续多无人机轨迹规划提供任务路径基础。

\subsection{考虑位置不确定性的多无人机轨迹规划方法}

针对多无人机在共享空域中的协同飞行问题，提出考虑位置不确定性的多无人机轨迹规划方法。该部分工作从编队飞行任务出发，将无人机和障碍物的位置不确定性引入碰撞建模过程，把传统的确定性几何避碰条件转化为可用于规划的概率安全约束。在此基础上，采用基于 MP 的轨迹规划方法，并将编队缩放纳入原语设计，使编队在保持整体构型的同时具备适应复杂环境的能力。通过在原语选择过程中综合考虑障碍物避让、机间避碰和位置不确定性影响，生成适用于多无人机协同飞行的安全轨迹。

\subsection{基于强化学习的多无人机协同运输控制方法}

针对重载条件下的运输执行问题，提出基于强化学习的多无人机协同运输控制方法。该部分工作围绕三机---负载系统的强耦合特性展开，将协同运输过程描述为共享回报的多智能体序贯决策问题，并在绳索约束下建立多机---负载耦合动力学模型。围绕这一任务模型，构建去中心化执行的多智能体强化学习控制框架，使策略输出能够通过底层控制接口转换为可执行控制量，从而实现多架无人机对重载负载的稳定协同运输。同时，为提高采样效率并缩短策略训练周期，进一步构建基于 MJX 的并行化仿真环境。该部分研究为重载条件下多无人机系统的稳定协同运输提供控制支撑，并与前述路径规划和轨迹规划内容相衔接，共同构成从全局路径生成、局部轨迹规划到协同运输控制的完整研究链路。

\section{章节安排}

全文共分为五章，各章内容安排如下。第一章为绪论，主要介绍低空物流中多无人机系统的研究背景与意义，梳理相关国内外研究现状，并给出研究内容、技术路线和章节安排。第二章为预备知识，主要介绍四旋翼建模基础、基于特殊欧氏群（Special Euclidean Group，SE(3)）的几何控制基础，以及强化学习与连续马尔可夫决策过程等相关理论，为后续章节的建模、规划与控制方法奠定基础。第三章研究能耗约束下基于改进关键可视图的路径规划方法，重点介绍复杂三维环境中的全局路径生成问题。第四章研究考虑位置不确定性的多无人机轨迹规划方法，给出满足概率安全约束和编队缩放要求的多无人机飞行轨迹生成方法。第五章研究基于强化学习的多无人机协同运输控制方法，构建面向重载运输任务的协同控制框架，并在并行化仿真环境中对所提方法进行验证与评估。

\section{本章小结}

本章首先介绍了低空物流中面向多无人机系统的研究背景及意义，然后围绕能耗约束下基于改进关键可视图的路径规划方法、考虑位置不确定性的多无人机轨迹规划方法以及基于强化学习的多无人机协同运输控制方法开展调研，归纳并分析了相关国内外研究现状，最后给出了研究内容、技术路线和章节安排。
